The gap isolated by Proposition~\ref{prop:v121-cofinal-rh} can be
viewed as a continuation and identification problem. In the screw-function
realization of Suzuki~\cite{Suzuki2026}, let $g$ be the Weil screw
function, $a=\log\lambda$, and $\lambda_a=\inf\sigma(A_a)$ as in
\eqref{eq:g2-monotone}. A certified window
$\mathsf W_\lambda\succeq0$ gives $\lambda_a\ge0$ and the associated
local positive kernel. The continuation theory of
Kre\u{\i}n--Langer~\cite{KreinLanger2014} distinguishes the existence of
positive extensions of local data from their uniqueness and from the
identification of a prescribed extension. Here the prescribed global
function is the arithmetic function $g$; Suzuki's global screw-function
criterion identifies its positivity with RH. A local certificate does
not establish that global identification. In particular we do not infer
that RH is equivalent to uniqueness of an extension of one fixed-window
kernel, or that uniqueness alone identifies the extension with $g$.
Any determinacy route must state and prove that additional identification.

Truncated moment problems and the Toeplitz--Hankel determinant theory
of Basor--Ehrhardt~\cite{BasorEhrhardt2002} provide a useful comparison,
not a further positivity theorem for this operator. The Fourier-basis
split in \texttt{evidence/diag\_circle\_split/} includes divided-difference
commutator terms; it is not an identification with the standard
single-symbol matrices $T(\phi)\pm H(\phi)$ covered by that reference.
Those numerical comparisons remain diagnostics. The cofinal estimate in
Proposition~\ref{prop:v121-cofinal-rh} is still a sufficient arithmetic
input of RH strength, not a consequence of extension terminology or a
finite list of positive sections. No G2 gap or RH claim follows.
